\paragraph{Optimizers as steepest descent under different norms.} SGD, Adam and Shampoo \citep{gupta2018shampoo} can be read as steepest descent under different norms \citep{bernstein2024old}. With its moving averages switched off, Adam becomes sign descent, steepest descent under the $\ell_\infty$ norm of the parameters; Shampoo without accumulation becomes spectral descent, the update $-\eta\,UV^\top$ for a gradient $G = USV^\top$, which is steepest descent under the spectral norm. Muon \citep{jordan2024muon} applies the same spectral-descent step to the momentum of each hidden weight matrix, approximating $UV^\top$ with a few Newton--Schulz iterations (Algorithm~\ref{alg:muon}), and leaves embeddings, the output head, gains and biases to AdamW. \citet{liu2025muon} scaled it to language-model pretraining. \citet{xie2025adam} argue that Adam's advantage over SGD comes from exploiting $\ell_\infty$ geometry: running Adam on a randomly rotated version of GPT-2's training loss made it converge much more slowly, worse even than a rotation-equivariant SGD variant, whereas SGD is rotation-equivariant and therefore unaffected. \citet{kunstner2024heavy} show that on language data with heavy-tailed class frequencies gradient descent makes slow progress on rare classes while Adam and sign descent are much less affected, and prove that in continuous time sign descent, unlike gradient descent, does not slow down on low-frequency classes. By analogy one might expect AdamW to have learned more of the weak input directions than SGD at a given error; at Garg scale SGD has learned slightly more (Section~\ref{sec:h1}).

\paragraph{Muon's implicit bias: parallel instead of sequential learning.} Several recent analyses ask which directions spectral optimizers learn first. \citet{vasudeva2025} study spectral gradient descent (the update $UV^\top$ of the gradient) on Gaussian-mixture data with linear and bilinear models: gradient descent learns the dominant principal components first, whereas spectral descent learns all of them at equal rates, which improves class-balanced loss early in training; depth amplifies the effect, and in experiments Muon and Shampoo help minority classes, especially early. \citet{wang2025muon} find that most of Muon's speed advantage over Adam in language-model training sits in the attention value/output weights and the MLPs, that Muon's updates have a more isotropic singular spectrum than Adam's, and prove that in a one-layer associative memory with imbalanced classes Muon learns all classes evenly while Adam's errors can be very uneven. \citet{dragutinovic2026} show that simplified Muon trades gradient descent's simplicity bias for speed. In deep linear networks it learns all modes at once; in one-layer multi-head linear attention on in-context regression with $x \sim \mathcal N(0,\Lambda)$ (one rank-one head per direction, small initialization), spectral gradient flow grows all heads at the same rate, so head $i$ finishes in time proportional to $\lambda_i^{-1/3}$ rather than one at a time as under gradient flow. In separate experiments (a spurious feature, a routing task that needs shared representations, and a two-layer softmax transformer on a ``cycle'' task) losing the simplicity bias hurt generalization. \citet{ma2026} prove that simplified Muon solves fixed-design in-context regression with a reparameterized one-layer linear transformer in a number of iterations independent of the condition number, while gradient descent and sign descent (Adam without momentum) slow down as it grows. The parallel learning is fragile, though: \citet{jacobs2026} show that small-batch gradient noise smooths the matrix-sign update toward one proportional to the gradient, restoring strong-modes-first learning in deep linear matrix factorization, with supporting experiments on vision transformers and language models. This matters for us because we train with batches of 64.

\paragraph{Partial flattening.} \citet{dong2026muonp} replace the singular values of the update by their $p$-th power, $US^pV^\top$, interpolating between Muon ($p=0$) and gradient descent ($p=1$). They show that this is steepest descent under the Schatten norm with exponent $1 + 1/p$, and find in language-model experiments that partial flattening helps fine-tuning while full flattening remains better for pretraining. \citet{shumaylov2026} find that, with tuned step sizes, GPT-2's validation loss is largely insensitive to the exact spectrum of the update: random singular values or a wide range of Schatten exponents match Muon, although SGD and an update that zeroes the largest singular values do not. Both studies judge spectra by loss or downstream accuracy. We use the same knob (Eq.~\ref{eq:knob}) to ask what the spectrum does to the learned algorithm at equal loss.

\paragraph{How this work differs.} The analyses of Muon's bias concern simplified Muon (no momentum, exact orthogonalization) in linear models, linear attention or associative memories, mostly in the gradient-flow or full-batch limit, and most judge optimizers by speed or loss. We train full Muon (momentum, Newton--Schulz) in multi-layer softmax transformers at Garg et al.'s scale with batch-64 noise; we compare optimizers at matched in-distribution error, so that speed is factored out and what remains is which algorithm has been learned; and we test the consequence out of distribution. The literature predicts effects of both signs there: learning weak directions earlier should help on rare inputs \citep{vasudeva2025,wang2025muon}, while losing gradient descent's simplicity bias can hurt generalization \citep{dragutinovic2026}. H2 predicts both, on different kinds of shifted queries.
